\documentclass[11pt]{article}

\usepackage{sectsty}
\usepackage{graphicx}
\usepackage{amsmath,amssymb}
% Margins
\topmargin=-0.45in
\evensidemargin=0in
\oddsidemargin=0in
\textwidth=6.5in
\textheight=9.0in
\headsep=0.25in

\title{A Review of \\ A Kernel Test for Causal Association via
	Noise Contrastive Backdoor Adjustment}
%\author{}
%\date{\today}

\begin{document}
	\maketitle	
	% Optional TOC
	% \tableofcontents
	% \pagebreak
	

This article proposes a non-parametric approach to testing for the presence of a causal effect by extending the work of Gretton et al. on the Hilbert-Schmidt Independence Criterion (HSIC) using g-computation principles. The specific contributions are as follows:
\begin{itemize}
	\item Introducing Backdoor-HSIC: Backdoor-HSIC is a novel extension of HSIC derived using importance-weighted covariance operators. The convergence properties of the corresponding estimators in causal settings have also been established. Also developed is a new optimization strategy to improve the effective sample size.
	\item Novel Permutation Strategy: This work provided a permutation strategy for backdoor-HSIC that ensures a permutation test with theoretically correct size for arbitrary treatment types.
	\item Empirical Evaluation: Demonstrating that backdoor-HSIC achieves correct size and exhibits strong power across various treatment types and with a large number of confounders when testing the do-null hypothesis (see the definition in the paper).
	\item Ablation Studies and Limitations: The article also conducted ablation studies to analyze backdoor-HSIC and characterized the circumstances under which it may become invalid.
\end{itemize}
	
\end{document}





